% ── Chapter 2: Related Work and Background Theory ─────────────
\chapter{Related Work and Background Theory}

\section{Antibiotic Susceptibility Testing}

The Kirby-Bauer disk diffusion assay was first systematically described and standardized by
Bauer, Kirby, Sherris, and Turck in 1966 and has since become the most widely adopted technique
for antimicrobial susceptibility testing in clinical microbiology laboratories worldwide
\cite{asm2009kirby,clsi2025}. Its longevity is attributable to its low equipment cost, procedural simplicity,
and adaptability to virtually any aerobically growing bacterial pathogen. The method depends on
the principle of antibiotic diffusion through semi-solid agar: when an antibiotic-impregnated
paper disk is placed on the surface of an inoculated agar plate, the drug molecules diffuse
radially outward from the disk, establishing an exponentially decreasing concentration gradient
as a function of distance. Where the local antibiotic concentration exceeds the minimum
inhibitory concentration (MIC) of the test organism, bacterial proliferation is suppressed,
producing a circular clearing zone on an otherwise confluent lawn of bacterial growth.

The diameter of this zone of inhibition (ZOI) is inversely related to the MIC: organisms with
high susceptibility (low MIC) exhibit large zones, while intrinsically or acquired-resistant
organisms exhibit small or absent zones. This relationship between ZOI diameter and MIC is not
linear but follows a mathematical relationship dependent on agar depth, disk potency, diffusion
coefficient, and incubation duration---all of which are carefully controlled by standardized
protocols. For clinical use, the continuous ZOI diameter measurement is converted to a categorical
susceptibility interpretation (Susceptible, Intermediate, or Resistant) by applying published
interpretive breakpoints that are themselves derived from clinical and microbiological data.
This method has been applied to a broad range of clinically relevant bacteria, including
\textit{Bacillus cereus} \cite{fitwi2020} and \textit{Bacillus subtilis} \cite{jiang2023},
among many other aerobic pathogens encountered in routine laboratory practice.

\subsection{Clinical Standards: CLSI and EUCAST}

Interpretive breakpoints---the threshold diameters used to classify organisms as Susceptible (S),
Intermediate (I), or Resistant (R)---are published and regularly updated by two principal
international standard-setting organizations: the Clinical and Laboratory Standards Institute
(CLSI) in the United States and the European Committee on Antimicrobial Susceptibility Testing
(EUCAST) in Europe \cite{eucast2023}. Although both organizations pursue the same clinical goal,
their methodological approaches to breakpoint derivation differ in important ways, and their
published thresholds for the same drug-organism combination can differ by several millimeters.

CLSI publishes its disk diffusion standards in the M02 document, currently in its fourteenth
edition \cite{clsi2025}, while susceptibility testing performance standards are also defined
in the M100 35th edition \cite{clsi_m100_2025}. Breakpoints are established by a
multidisciplinary working group that integrates MIC distribution data from large collections
of clinical isolates, PK/PD modeling of drug exposure at standard therapeutic doses, and
clinical outcome data from randomized controlled trials. The Intermediate (I) category in
the CLSI system traditionally indicated a buffer zone where technical uncertainty might lead
to misclassification, but more recent editions have increasingly aligned with EUCAST's concept
of the ``Susceptible, Increased Exposure'' (I) category, which indicates susceptibility when
drug exposure is maximized through dosing optimization.

EUCAST, which publishes its breakpoint tables annually with version numbers (the current version
referenced in this project is 13.1 \cite{eucast2023}), derives its breakpoints from
pharmacokinetic/pharmacodynamic (PK/PD) targets, epidemiological cutoff values (ECOFFs), and
clinical outcome data \cite{giske2022}. ECOFFs represent the upper end of the MIC distribution
for wild-type bacterial populations without acquired resistance mechanisms; any organism with an
MIC above the ECOFF is defined as non-wild-type (NWT) and presumed to harbor resistance
determinants. EUCAST breakpoints are frequently more aggressive in separating susceptible from
resistant populations, reflecting a European pharmacological framework that differs from the
United States standard \cite{kahlmeter2019}. Both sets of breakpoints are implemented in the
database of the present system, allowing laboratories to select the applicable standard
according to their regulatory and clinical context.

\subsection{Clinical Significance of Measurement Precision}

The clinical consequence of measurement error in disk diffusion is direct and potentially severe.
Because breakpoints are expressed as integer millimeter thresholds, a single millimeter of
measurement error can shift a classification from Susceptible to Resistant or vice versa. This
misclassification has two dangerous modes: false susceptibility (classifying a resistant organism
as susceptible) leads to treatment failure and potentially fatal outcome; false resistance
(classifying a susceptible organism as resistant) leads to unnecessary escalation to
second-line or third-line agents, contributing to resistance selection pressure and higher
costs. Humphries et al.\ \cite{humphries2018} defined a mean absolute error of less than
3\,mm as the boundary of clinical acceptability for automated AST measurement tools, and this
threshold serves as the primary performance target for the present project.

\section{Deep Learning for Object Detection}

Object detection in computer vision is the task of simultaneously localizing and classifying
one or more objects within a single image. Machine learning, broadly defined, encompasses
supervised, unsupervised, and reinforcement learning paradigms \cite{alam2023ml,murphy2025rl},
of which supervised deep learning has proven most effective for visual recognition tasks.
Convolutional Neural Networks (CNNs) form the backbone of most modern object detectors
\cite{purwono2023cnn}, and systematic reviews of medical deep learning applications have
confirmed their suitability for a wide range of clinical image analysis tasks \cite{egger2022}.
A detector must output, for each detected object, both a bounding box specifying the object's
spatial extent and a class label identifying its category. This is fundamentally more complex
than image classification, which assigns a single label to an entire image, because the number
and location of objects are unknown a priori and must be inferred from the image content.
Deep learning approaches have dominated object detection benchmarks since the publication of
the Region-based CNN (R-CNN) family and YOLO models, progressively improving both accuracy
and inference speed.

Modern deep learning object detectors are generally categorized into two architectural families:
two-stage detectors, which first generate region proposals and then classify each proposal; and
one-stage (single-shot) detectors, which predict bounding boxes and class probabilities directly
in a single forward pass over the full image. Each family offers distinct trade-offs between
detection accuracy, computational cost, and inference latency.

\subsection{Faster R-CNN}

Faster R-CNN (Region-based Convolutional Neural Network) was introduced by Ren et al.\ in 2015
as a landmark advance that unified region proposal generation with the detection head into a
single end-to-end trainable network \cite{ren2015faster}. Prior to Faster R-CNN, the R-CNN
and Fast R-CNN pipelines required a separate, computationally expensive selective search step
to generate region proposals, which bottlenecked inference speed. Faster R-CNN replaced this
with a Region Proposal Network (RPN) that shares the convolutional backbone with the detection
head, enabling near real-time proposal generation.

\begin{figure}[H]
  \centering
  \safeimg[width=0.82\textwidth]{figures/faster_rcnn_architecture.png}{Faster R-CNN two-stage detection architecture}
  \caption{Architecture of Faster R-CNN, showing the shared convolutional backbone, the Region
    Proposal Network (RPN) producing candidate object regions, and the detection head that
    classifies and refines each proposal.}
  \label{fig:faster_rcnn}
\end{figure}

The Faster R-CNN architecture operates in two stages. In the first stage, the backbone
convolutional network (in this project, ResNet-50 \cite{he2016deep}) extracts a rich
spatial feature map from the full input image. The RPN then slides a small network over
this feature map and, at each spatial location, predicts objectness scores and bounding
box regression offsets for a set of predefined anchor boxes of multiple scales and aspect
ratios. Anchor boxes that score above an objectness threshold are retained as region proposals
(typically 256--512 per image after non-maximum suppression). In the second stage, the
RoI Pooling (or RoI Align) operation extracts a fixed-size feature vector for each proposal
from the shared backbone feature map, and a fully connected classification head assigns a
class label and a refined bounding box to each region.

The use of ResNet-50 as the feature extraction backbone is particularly well-suited to this
project. He et al.\ demonstrated that residual connections---shortcut connections that bypass
one or more convolutional layers---enable training of significantly deeper networks without
the gradient degradation problems that had previously limited depth \cite{he2016deep}. ResNet-50
comprises 50 layers organized into residual blocks, and its ImageNet pretraining provides a
rich initialization of low-level edge, texture, and mid-level shape features that transfer
effectively to the AST domain even with a relatively small number of labeled AST images.
The ResNet-50 backbone in this project is initialized with ImageNet-pretrained weights and
fine-tuned end-to-end on the AST dataset.

Faster R-CNN is recognized for its high detection accuracy, particularly in complex scenes
containing overlapping or densely packed objects---a characteristic directly relevant to AST
plates where zones from adjacent disks may partially overlap. Its two-stage architecture
allows the classification head to operate on high-quality, spatially refined proposals rather
than on a dense grid, reducing false positive rates. The primary limitation of Faster R-CNN
is inference speed: the two-stage pipeline introduces latency that is acceptable for offline
analysis but may limit real-time processing throughput.

\subsection{YOLOv11}

YOLO (You Only Look Once) models are single-stage detectors that reformulate object detection
as a unified regression problem. Rather than generating region proposals and classifying them
separately, YOLO models predict bounding box coordinates and class probabilities directly from
the full image in a single forward pass of the network \cite{wang2024yolov11,khanam2024yolov11}.
This design eliminates the region proposal stage entirely, enabling dramatically faster inference
at the cost of some accuracy on small or densely overlapping objects in early versions.
Comparative studies have confirmed YOLOv11's competitive performance relative to prior YOLO
generations and Faster R-CNN across multiple detection tasks \cite{sharma2024yolo_comparison}.

\begin{figure}[H]
  \centering
  \safeimg[width=0.82\textwidth]{figures/yolov11_architecture.png}{YOLOv11 single-stage detection architecture}
  \caption{Schematic architecture of YOLOv11, illustrating the backbone feature extraction
    network, the feature pyramid neck for multi-scale detection, and the decoupled detection
    heads for bounding box regression and class prediction.}
  \label{fig:yolov11}
\end{figure}

YOLOv11, the latest version in the YOLO family at the time of this project and developed by
Ultralytics, incorporates several architectural advances over its predecessors \cite{wang2024yolov11}.
The backbone employs enhanced C3k2 blocks (a reparameterized convolutional design) for improved
feature extraction efficiency. The neck uses a Path Aggregation Feature Pyramid Network
(PA-FPN) to fuse features from multiple backbone scales, enabling detection of objects across
a wide range of sizes---from small antibiotic disks to large inhibition zones that may span
a significant fraction of the plate. The detection head is decoupled, with separate branches
predicting bounding box regression and class probabilities independently, which has been shown
to improve convergence and accuracy in one-stage detectors.

YOLOv11 is offered in five size variants (n, s, m, l, x) that provide a continuous trade-off
between model capacity and inference speed. In this project, the nano (n) variant, YOLOv11n,
was selected to prioritize inference speed and suitability for eventual deployment on
resource-constrained hardware, such as laboratory workstations without dedicated GPUs or
on-device edge hardware. Despite being the smallest variant, YOLOv11n achieves competitive
detection accuracy on the AST dataset, benefiting from pretraining on the large-scale COCO
benchmark and fine-tuning on the project's annotated AST images. The model detects two
object classes: \texttt{antibiotic} (the drug-impregnated disk, used as an OCR target and
calibration reference) and \texttt{disk-zone} (the surrounding inhibition zone, from which
the diameter measurement is extracted).

\subsection{Hough Circle Transform}

The Hough Circle Transform is a classical computer vision technique for detecting circular
structures in images without requiring any learned model or labeled training data. It operates
by transforming the image from the spatial domain to a parameter space defined by three
dimensions: the $x$ and $y$ coordinates of the circle center and the circle radius $r$.
For each edge point detected in the image (typically via Canny edge detection), the algorithm
votes for all circles in parameter space that could have produced that edge point. Circles that
accumulate votes above a threshold in the parameter-space accumulator correspond to genuine
circles in the image. Alternative circle detection approaches based on learning automata
\cite{cuevas2014}, deep learning \cite{ercan2020}, and gradient-based variants \cite{shehata2015hough}
have been proposed to address the sensitivity limitations of the classical transform under real-world imaging conditions.

\begin{figure}[H]
  \centering
  \safeimg[width=0.75\textwidth]{figures/hough_circle_diagram.png}{Hough Circle Transform parameter space voting diagram}
  \caption{Illustration of the Hough Circle Transform. Each edge point votes for a cone of
    candidate circles in the $(x_c, y_c, r)$ parameter space; genuine circles produce
    prominent peaks in the accumulator.}
  \label{fig:hough}
\end{figure}

The Hough Circle Transform is computationally efficient and produces no false positives for
clearly circular zones under controlled imaging conditions. However, it is sensitive to several
sources of degradation common in real AST images. Uneven agar surface reflectance or non-uniform
lighting creates spurious edge responses that generate noise in the accumulator. Zones that
are not perfectly circular---due to zone overlap, agar surface irregularities, or atypical
bacterial growth patterns---do not accumulate votes cleanly, leading to missed detections or
inaccurate radius estimates. Furthermore, the transform requires careful hand-tuning of multiple
parameters (the inverse resolution ratio \texttt{dp}, the minimum inter-center distance
\texttt{minDist}, and the Canny edge thresholds \texttt{param1} and \texttt{param2}) that
interact in complex ways and whose optimal values vary across images with different scales,
lighting conditions, and bacterial species. In this project, Hough Circle Transform parameters
were tuned via grid search on the full dataset, representing a best-case scenario for this
method. It serves as a principled classical baseline against which the learned models are evaluated.

\section{Optical Character Recognition for Antibiotic Disk Labeling}

Antibiotic disks carry short alphanumeric codes printed on their surfaces, indicating the
antibiotic name and its concentration (for example, ``OX 1'' for oxacillin 1\,$\mu$g, or
``AMX 25'' for amoxicillin 25\,$\mu$g). Optical Character Recognition (OCR), whose development
history spans several decades of engineering progress \cite{wang2023ocr_history}, provides the
computational means to read these labels automatically from photographic images. Correctly
reading these labels is essential for associating each detected ZOI with the correct antibiotic,
which in turn enables the correct breakpoint to be retrieved from the CLSI/EUCAST database.
This is a challenging OCR task for several reasons: the printed text is small (typically 2--4 mm
in height), the contrast between ink and disk surface is variable across disk brands, the disk
surface is slightly curved, images may be taken at sub-optimal angles, and multiple disks may
be partially occluded.

\subsection{EasyOCR and the CRAFT Architecture}

EasyOCR \cite{jaidedai2024easyocr} is an open-source, deep learning-based OCR engine
supporting over 80 languages and scripts. Its text detection stage is based on the CRAFT
(Character Region Awareness for Text Detection) model introduced by Baek et al.\
\cite{baek2019craft}. CRAFT detects individual character regions and the affinity between
adjacent characters rather than treating word-level bounding boxes as the atomic detection
unit. This character-level detection approach is particularly well-suited to the antibiotic
labeling context, where short strings of alphanumeric characters must be localized precisely
on small, textured disk surfaces.

The CRAFT network produces two output maps: a region score map (indicating the probability
that each pixel belongs to the center of a character) and an affinity score map (indicating
the probability that two adjacent characters belong to the same word). These two maps are
combined to segment individual characters and group them into word-level text instances.
The detected character regions are then passed to a sequence recognition network---a
Convolutional Recurrent Neural Network (CRNN)---that transcribes the character sequence
into text using a connectionist temporal classification (CTC) decoder.

EasyOCR \cite{jaidedai2024easyocr} was compared against TyphoonOCR \cite{typhoonocr2025},
a Thai-focused OCR model developed by OpenTyphoon.ai, for the antibiotic label reading task.
Although TyphoonOCR offers strong performance on Thai-script text, its performance on Latin
alphanumeric strings under the challenging conditions of disk photography was inferior to
EasyOCR \cite{baek2019craft,jaidedai2024easyocr}. EasyOCR was therefore selected as the
OCR engine for the final pipeline.

\section{Image Preprocessing Techniques}

Image preprocessing is a critical enabler of both the object detection models and the OCR
pipeline, directly affecting the signal-to-noise ratio of the features extracted by downstream
algorithms \cite{qi2021image_enhancement}. Image segmentation techniques \cite{jasim2021segmentation,dhanachandra2015}
and document binarization methods \cite{yang2024binarization} inform the preprocessing choices
made in this project to address the specific visual degradations typical of laboratory AST photography.

\begin{itemize}
  \item \textbf{Global thresholding (Otsu's method):} Otsu's method automatically determines
        an optimal intensity threshold for converting a grayscale image to a binary (black and
        white) representation by minimizing intra-class variance between the foreground (text)
        and background pixel populations. This is applied to disk crops prior to OCR to separate
        ink characters from the disk surface background, improving character segmentation quality.

  \item \textbf{Contrast Limited Adaptive Histogram Equalization (CLAHE):} Unlike global
        histogram equalization, which computes a single mapping function for the entire image,
        CLAHE applies histogram equalization independently to small tiles of the image (with a
        configurable tile size and contrast limit) and interpolates smoothly between them.
        This normalizes local contrast across regions of the image that differ substantially
        in mean illumination---a common occurrence in laboratory photographs where the agar plate
        center is brighter than the periphery due to overhead lighting geometry. CLAHE is applied
        to disk crops as a preprocessing step before OCR, improving the visibility of low-contrast
        text.

  \item \textbf{Gaussian Blur:} A Gaussian low-pass filter is applied to the full plate image
        before edge detection (for the Hough Transform) and before thresholding, attenuating
        high-frequency sensor noise that would otherwise produce spurious edge responses.
        The blur kernel size is selected to preserve the relevant scale of zone boundaries
        while suppressing noise.

  \item \textbf{Morphological operations (dilation and erosion):} Morphological dilation
        expands foreground (text) regions, while erosion shrinks them. Applied in sequence
        (as opening or closing operations), they can fill small gaps in characters or remove
        small noise patches from the binary image. In the OCR preprocessing pipeline,
        the choice and extent of morphological operations are applied adaptively based on the
        estimated text pixel ratio in the binarized crop: images with very sparse text pixels
        (thin characters) are dilated to restore stroke connectivity, while images with
        excessively dense pixels (thick strokes or merged characters) are eroded.

  \item \textbf{Image rotation correction:} The orientation of antibiotic disks on the plate
        is not controlled during laboratory setup---disks may be placed at any angle. As a
        result, the printed label on a disk may appear rotated at an arbitrary angle relative
        to the image frame. To handle this, the OCR pipeline applies EasyOCR at four canonical
        orientations (0$^\circ$, 90$^\circ$, 180$^\circ$, and 270$^\circ$) and selects the
        highest-confidence recognition result across all four orientations.
\end{itemize}

\section{Related Work in Automated AST Analysis}

\subsection{Deep Learning Approaches for AST Image Analysis}

The application of deep learning and computer vision to AST image analysis has been an active
area of research in the past decade, driven by the same combination of clinical need and
algorithmic progress that motivates this project. The broader role of AI and machine learning
in predicting and combating antimicrobial resistance has been comprehensively reviewed by
Bilal et al.\ \cite{bilal2025}, while Liao et al.\ \cite{liao2025} have surveyed advancements
specifically in AI-driven drug sensitivity testing research.
Pascucci et al.\ \cite{pascucci2021,pascucci2020preprint} developed the first AI-based mobile
application for automated antibiotic resistance testing, demonstrating that smartphone cameras
combined with deep learning could produce clinically acceptable ZOI measurements in
resource-limited settings.
Rajasekar et al.\ \cite{rajasekar2022} introduced Ab.ai, a specialized AI tool for reporting
antibiograms automatically, achieving high accuracy on disk diffusion plates in standardized
laboratory conditions.
Yu et al.\ \cite{yu2025} demonstrated an AI-driven hue-contrast method for inhibition zone
identification, providing an alternative to bounding-box regression.
Callebaut et al.\ \cite{callebaut2024} evaluated the Radian\textregistered{} in-line carousel
system for automated disk diffusion testing, establishing performance benchmarks
against expert manual reading that this project adopts as comparison targets.

Rajasekar et al.\ \cite{rajasekar2023} conducted a comprehensive
systematic review of published AI-based image analysis methods for AST and found that deep
learning models consistently and substantially outperform classical image processing approaches---
including Hough Transform, active contours, and template matching---particularly in the
scenarios most relevant to clinical practice: plates with dense disk arrays, overlapping zones,
and heterogeneous bacterial growth. The review surveyed methods ranging
from traditional convolutional neural networks applied to cropped disk regions to full end-to-end
detection pipelines similar to the one developed in this project.

Among the key findings of Rajasekar et al.\ were that the greatest performance gains from deep
learning occur specifically in the handling of overlapping zones---a scenario that fundamentally
breaks the assumption of isolated circular shapes underlying the Hough Transform approach---and
in robustness to photographic variation such as non-uniform illumination, perspective distortion,
and image compression artifacts \cite{rajasekar2023}. These findings validate the architectural
choices of the present project: using deep learning object detection as the primary approach
and relegating Hough Transform to the role of a classical baseline rather than a primary system
component. The review also noted that the relatively small scale of most published AST datasets
(typically below 500 images) has been a limiting factor for deep learning performance, motivating
the present project's investment in synthetic data generation to reach a training set of over
2,000 images.

\subsection{Evaluation Standards for Automated AST Methods}

Humphries et al.\ \cite{humphries2018} addressed the methodological question of how to properly
evaluate an automated AST measurement system relative to the gold standard of manual reading by
a trained microbiologist. Their best-practices document, published by the CLSI Methods
Development and Standardization Working Group, defined several important evaluation criteria
that the present project adopts. Most critically for this work, they established that a mean
absolute error (MAE) of less than 3\,mm between automated and manual ZOI diameter measurements
represents the threshold of clinical acceptability---that is, below this error level, the
systematic measurement differences are unlikely to produce a different S/I/R classification
for the majority of drug-organism breakpoint combinations currently in use.

Humphries et al.\ also emphasized the importance of evaluating automated methods across a
diverse collection of bacterial species, antibiotic classes, and plate configurations, since
errors that are insignificant for one drug-organism combination may be clinically decisive
for another \cite{humphries2018}. This guidance informs the evaluation design of the present
project, which reports MAE across the full held-out test set of 235 images drawn from the
real plate collection. Additionally, Humphries et al.\ noted that systematic (directional)
bias is as important to characterize as random error, since a consistent over- or under-estimation
of zone diameter can be corrected by a post-processing bias correction step---which is precisely
what the calibration refinement stage of the present system implements.
